\section{Reductions and Newton potentials}\label{sec:preliminaries}

We begin with consequences of positivity that do not require the
subcritical inequality.  Throughout this section, a solution means a pair
\[
 u,v\in C^2(\R^n\setminus\{0\})\cap C(\R^n),\qquad u,v>0
 \quad\hbox{on }\R^n,
\]
where $n\ge2$, $p,q>0$, and $A,B\in\R$, satisfying
\[
 -\Delta u=|x|^A v^p,\qquad -\Delta v=|x|^B u^q
 \quad\hbox{on }\R^n\setminus\{0\}.
\]
We write $B_R(z)=\{x:|x-z|<R\}$ and $B_R=B_R(0)$.

\begin{lemma}[Origin reductions]\label{lem:reductions}
The existence of a solution implies $n\ge3$ and $A,B>-2$.
\end{lemma}

\begin{proof}
For a function $h$ continuous on each sphere about the origin, write
\[
 \overline h(r)=\frac1{|\Sph|}\int_{\Sph}h(r\theta)\,\dd\theta.
\]
Integration of the equations on annuli gives, for $r>0$,
\begin{equation}\label{eq:spherical-flux}
 \begin{split}
  \bigl(r^{n-1}\overline u'(r)\bigr)'
     &=-r^{n-1+A}\overline{v^p}(r),\\
  \bigl(r^{n-1}\overline v'(r)\bigr)'
     &=-r^{n-1+B}\overline{u^q}(r).
 \end{split}
\end{equation}
In particular, each flux on the left is strictly decreasing.  Its extended
limit as $r\downarrow0$ exists.  If that limit were nonzero or infinite,
the flux would have a fixed sign and absolute value bounded below by a
positive constant for all sufficiently small $r$.  Since
$\int_0^1 r^{1-n}\,\dd r=\infty$, integration would contradict boundedness
of the corresponding spherical mean at zero.  Thus both limits are zero,
and both fluxes are strictly negative for every $r>0$.

If $n=2$, then $r\overline u'(r)\le\overline u'(1)<0$ for $r\ge1$.
Consequently
$\overline u(r)\le\overline u(1)+\overline u'(1)\log r$, contradicting
positivity for large $r$.

Continuity and $v(0)>0$ give $v^p\ge c>0$ near zero.  Integrating the first
equation in \eqref{eq:spherical-flux} over $[r/2,r]$ and using the negative
flux at $r/2$, we obtain
\[
 -\overline u'(r)
 \ge c r^{1-n}\int_{r/2}^r s^{n-1+A}\,\dd s
 =c_A r^{1+A},\qquad c_A>0.
\]
The equality remains valid when $n+A=0$, with the corresponding constant
containing $\log2$.  If $A\le-2$, integrating this inequality towards
zero contradicts continuity of $u$ there.  Interchanging the two equations
proves $B>-2$.
\end{proof}

Henceforth assume $n\ge3$ and $A,B>-2$.  Define, off zero,
\begin{equation}\label{eq:sources-kernel}
 \begin{gathered}
 f(x)=|x|^B u(x)^q,\qquad g(x)=|x|^A v(x)^p,\\
 K(x)=c_n|x|^{2-n},\qquad c_n=\frac1{(n-2)|\Sph|}.
 \end{gathered}
\end{equation}
We set $f(0)=g(0)=0$, which has no effect on any integral.  All unspecified
integrals below are Lebesgue integrals.  A constant called \emph{universal}
may depend on $n,p,q,A,B$ and on explicitly indicated auxiliary parameters,
but not on the particular solution.  In contrast, constants denoted by
$C_{u,v}$ are allowed to depend on the solution.

\begin{lemma}[Distributional extension and potentials]\label{lem:potential}
For every solution with $n\ge3$ and $A,B>-2$, the sources $f,g$ belong to
$L^1_{\mathrm{loc}}(\R^n)$ and
\[
 -\Delta u=g,\qquad -\Delta v=f
 \quad\hbox{in distributions on }\R^n.
\]
The Newton potentials are finite everywhere and satisfy
\begin{equation}\label{eq:potential-lower}
 K*g\le u,\qquad K*f\le v \quad\hbox{on }\R^n.
\end{equation}
There are constants $c_u,c_v\ge0$ such that
\begin{equation}\label{eq:potential-constants}
 u=c_u+K*g,\qquad v=c_v+K*f
\end{equation}
almost everywhere on $\R^n$ and pointwise on $\R^n\setminus\{0\}$.
On the latter set the potentials are continuously differentiable, and
\begin{equation}\label{eq:potential-gradients}
 \begin{split}
 \nabla u(x)&=\int_{\R^n}\nabla K(x-y)g(y)\,\dd y,\\
 \nabla v(x)&=\int_{\R^n}\nabla K(x-y)f(y)\,\dd y,\\
 \nabla K(z)&=-(n-2)\frac{z}{|z|^2}K(z).
 \end{split}
\end{equation}
Both gradient integrals converge absolutely at every $x\ne0$.
\end{lemma}

\begin{proof}
Local boundedness of $u,v$ gives
$g(x)\le C_{u,v}|x|^A$ and $f(x)\le C_{u,v}|x|^B$ near zero.
Since $n+A,n+B>0$, both sources are locally integrable.  Let
$\varphi\in C_c^\infty(\R^n)$ and choose
$\xi\in C_c^\infty(\R^n)$ equal to one near zero.  Test the equation for
$u$ away from zero against
$\varphi(x)(1-\xi(x/\rho))$.  After expanding its Laplacian, the terms
involving first and second derivatives of $\xi(x/\rho)$ have absolute
values bounded respectively by $C_{u,v,\varphi,\xi}\rho^{n-1}$ and
$C_{u,v,\varphi,\xi}\rho^{n-2}$.  They tend to zero.  The remaining
terms pass to the limit by local integrability, proving the first
distributional equation.  The second follows in the same way.

Let $\zeta_\rho$ be a nonnegative, compactly supported smooth approximate
identity, and put $u_\rho=u*\zeta_\rho$ and $g_\rho=g*\zeta_\rho$.
Then $u_\rho>0$ and $-\Delta u_\rho=g_\rho\ge0$ on all of $\R^n$.
Two integrations of the spherical flux formula, now about an arbitrary
point $x$, give
\[
 u_\rho(x)-\frac1{|\Sph|}\int_{\Sph}u_\rho(x+R\theta)\,\dd\theta
 =c_n\int_{|x-y|<R}
       \bigl(|x-y|^{2-n}-R^{2-n}\bigr)g_\rho(y)\,\dd y.
\]
The spherical average is positive.  Monotone convergence as $R\to\infty$
therefore yields $K*g_\rho\le u_\rho$.  As $\rho\downarrow0$,
$u_\rho(x)\to u(x)$ for every $x$, while $g_\rho\to g$ almost everywhere.
Fatou's Lemma proves $K*g(x)\le u(x)$ at every point.  The same argument
proves $K*f\le v$.

The normalization of $K$ gives $-\Delta K=\delta_0$ in distributions:
$K$ is harmonic off zero,
$-\int_{\partial B_r}\partial_rK=1$, and
$\int_{\partial B_r}K=O(r)$, so integration by parts outside $B_r$
and passage to $r=0$ gives the assertion.  To apply this to $K*g$, note
that for every compactly supported smooth $\varphi$,
\[
 \int_{\R^n}g(y)\int_{\R^n}K(x-y)|\Delta\varphi(x)|\,\dd x\,\dd y
 =\int_{\R^n}(K*g)(x)|\Delta\varphi(x)|\,\dd x
 \le\int_{\R^n}u(x)|\Delta\varphi(x)|\,\dd x<\infty.
\]
Fubini's Theorem thus gives $-\Delta(K*g)=g$.  The difference
$h=u-K*g$ is a nonnegative locally integrable harmonic distribution.
Every $h*\zeta_\rho$ is a nonnegative entire harmonic function.
Such a function $H$ is constant: for $b=|x-z|$, the mean-value property
and $B_R(x)\subset B_{R+b}(z)$ imply
\[
 H(x)\le\left(\frac{R+b}{R}\right)^n H(z).
\]
Letting $R\to\infty$ and then interchanging $x,z$ proves equality.
Passing back to distributions shows that $h$ itself is a nonnegative
constant.  The argument for $v-K*f$ is identical.

Fix $x_0\ne0$.  The sources are locally $C^1$ near $x_0$, because $u,v$
are strictly positive there and $p,q>0$.  Split $g$ by a smooth cutoff
supported near $x_0$ and equal to one on a smaller neighborhood.  For
the compactly supported part, local integrability of $K$ and $\nabla K$
gives a continuously differentiable convolution with the asserted
derivative.  For the other part, $|x-y|$ stays bounded away from zero
as $x$ varies in a sufficiently small neighborhood of $x_0$.  There
both $K(x-y)$ and $|\nabla K(x-y)|$ are bounded by a fixed multiple of
$K(x_0-y)$.  The latter is integrable against $g$ by
\eqref{eq:potential-lower}.  Dominated convergence proves continuity
and differentiation of this part as well.  The almost-everywhere identity
in \eqref{eq:potential-constants} is therefore pointwise off zero, and
\eqref{eq:potential-gradients} follows.  The proof for $f$ is the same.
\end{proof}

For $x,y\in B_R$, the kernel satisfies $K(x-y)\ge cR^{2-n}$.  Thus
Lemma~\ref{lem:potential} gives the centered lower bounds
\begin{equation}\label{eq:ball-lower}
 u(x)\ge cR^{2-n}\int_{B_R}g,\qquad
 v(x)\ge cR^{2-n}\int_{B_R}f\qquad(x\in B_R),
\end{equation}
where $c>0$ depends only on $n$.

\begin{lemma}[Exclusion of small products]\label{lem:low-exponents}
A solution with $n\ge3$ and $A,B>-2$ necessarily has $pq>1$.
\end{lemma}

\begin{proof}
Set $m_u(R)=\inf_{B_R}u$ and $m_v(R)=\inf_{B_R}v$.  Both are positive
by positivity and continuity on $\overline{B_R}$.  Integrating the
weights in \eqref{eq:ball-lower} yields
\[
 m_u(R)\ge cR^{2+A}m_v(R)^p,\qquad
 m_v(R)\ge cR^{2+B}m_u(R)^q.
\]
Combining them and dividing by $m_u(R)^{pq}>0$ gives
\[
 m_u(R)^{1-pq}\ge cR^{2+A+p(2+B)}.
\]
The exponent on the right is positive.  If $pq<1$, the left side is
at most $u(0)^{1-pq}$, and if $pq=1$, it is one.  Either case contradicts
the inequality as $R\to\infty$.
\end{proof}

We may consequently assume
$n\ge3$, $A,B>-2$, and $pq>1$.  Introduce
\begin{equation}\label{eq:parameters}
 \begin{gathered}
 \alpha=\frac{2+A+p(2+B)}{pq-1},\qquad
 \beta=\frac{2+B+q(2+A)}{pq-1},\\
 a=\frac1{q+1},\qquad b=\frac1{p+1}.
 \end{gathered}
\end{equation}
Both $\alpha$ and $\beta$ are positive, and
\[
 p\beta=\alpha+2+A,\qquad q\alpha=\beta+2+B.
\]
For every $R>0$, the pair
\begin{equation}\label{eq:scaling}
 \begin{split}
 u_R(x)&=R^\alpha u(Rx),\qquad v_R(x)=R^\beta v(Rx),\\
 f_R(x)&=|x|^B u_R(x)^q=R^{\beta+2}f(Rx),\\
 g_R(x)&=|x|^A v_R(x)^p=R^{\alpha+2}g(Rx)
 \end{split}
\end{equation}
again solves the same system, with the same parameters and regularity.
These are dilations about the origin, where the weights are centered.

\begin{proposition}[Universal centered estimates]\label{prop:uniform}
Fix $n\ge3$, $p,q>0$, $A,B>-2$, and $pq>1$.  There is a universal
constant $C$ such that every solution satisfies, for all $R>0$,
\begin{equation}\label{eq:source-bounds}
 \int_{B_R}f\le CR^{n-2-\beta},\qquad
 \int_{B_R}g\le CR^{n-2-\alpha}.
\end{equation}
The constants in \eqref{eq:potential-constants} vanish, so that
\begin{equation}\label{eq:potentials}
 u=K*g,\qquad v=K*f\quad\hbox{on }\R^n\setminus\{0\}.
\end{equation}
Moreover,
\begin{equation}\label{eq:uniform}
 \begin{split}
 &\int_{B_4}(f+g+u+v)\le C,\\
 &\int_{|x-y|\ge\ell}K(x-y)(f(y)+g(y))\,\dd y
       \le C\ell^{2-n}
       \qquad (x\in B_2,\ 0<\ell\le1).
 \end{split}
\end{equation}
\end{proposition}

\begin{proof}
Write $M_f=\int_{B_1}f$ and $M_g=\int_{B_1}g$.  The lower bounds
\eqref{eq:ball-lower}, followed by integration of $|x|^B$ and $|x|^A$
on $B_1$, give
\[
 M_f\ge c_fM_g^q,\qquad M_g\ge c_gM_f^p,
\]
where $c_f,c_g>0$ are universal.  In particular,
$M_f\ge c_fc_g^qM_f^{pq}$, and $M_f>0$ and $pq>1$ give a universal
upper bound for $M_f$.  The interchanged argument bounds $M_g$.
Apply these bounds to the scaled pair \eqref{eq:scaling}; changing
variables proves \eqref{eq:source-bounds}.  The argument uses only
positive powers, and requires neither $p\ge1$ nor $q\ge1$.

If $c_u>0$ in \eqref{eq:potential-constants}, then $u\ge c_u$ almost
everywhere and
\[
 \int_{B_R}f\ge c_u^q\int_{B_R}|x|^B\,\dd x
 =c_u^q\frac{|\Sph|}{n+B}R^{n+B}.
\]
This contradicts \eqref{eq:source-bounds} as $R\to\infty$, since
$n+B-(n-2-\beta)=B+2+\beta=q\alpha>0$.
Likewise $c_v=0$, proving \eqref{eq:potentials}.

For $x\in B_4$ and $j\ge3$, the annulus
$2^j\le|y|<2^{j+1}$ has $|x-y|\ge2^{j-1}$.
Its contribution to the $f$ potential is at most
\[
 C2^{j(2-n)}\int_{B_{2^{j+1}}}f\le C2^{-j\beta},
\]
and its contribution to the $g$ potential is at most $C2^{-j\alpha}$.
The two series converge because $\alpha,\beta>0$.  Thus the potentials
of the sources restricted to $\R^n\setminus B_8$ are uniformly bounded
on $B_4$.  On the other hand,
\[
 \sup_{y\in B_8}\int_{B_4}K(x-y)\,\dd x
 \le\int_{B_{12}}K(z)\,\dd z<\infty.
\]
Tonelli's Theorem, the source bounds on $B_8$, and
\eqref{eq:potentials} now bound $\int_{B_4}(u+v)$.  The source part
of the first estimate in \eqref{eq:uniform} follows directly from
\eqref{eq:source-bounds}.

Finally, on $|x-y|\ge\ell$ the kernel is at most $c_n\ell^{2-n}$.
This bounds the contribution from $y\in B_8$ by $C\ell^{2-n}$.
The contribution from $y\notin B_8$ has already been bounded by $C$;
since $0<\ell\le1$ and $n\ge3$, this proves the second estimate in
\eqref{eq:uniform}.
\end{proof}

The bounds \eqref{eq:source-bounds} are centered at the origin.  For
example, when a later argument uses a ball $B_\ell(x)\subset B_4$, its
source mass is bounded by the mass of $B_4$ in \eqref{eq:uniform}.
No translation of the weighted equation or source estimate is needed.
Also, local energy is finite for each solution:
\[
 fu=|x|^Bu^{q+1},\qquad gv=|x|^Av^{p+1}
 \quad\hbox{belong to }L^1_{\mathrm{loc}}(\R^n).
\]
This statement follows from local boundedness and $A,B>-2$; it is not
yet a bound uniform over solutions.  We assume no finiteness of either
energy integral on all of $\R^n$.

To connect the centered dilation with the subcritical hypothesis, define
\begin{equation}\label{eq:gap}
 \begin{split}
 \gamma
 &:=a(n+B)+b(n+A)-(n-2)\\
 &=(1-a-b)(\alpha+\beta+2-n),\\
 1-a-b&=\frac{pq-1}{(p+1)(q+1)}>0.
 \end{split}
\end{equation}
Indeed, the numerator of $\alpha+\beta$ is
$(q+1)(2+A)+(p+1)(2+B)$, and substitution gives the identity.
The strict inequality in Theorem~\ref{thm:main} is precisely
$\gamma>0$, and hence implies $\alpha+\beta+2-n>0$.
Finally, changing variables in \eqref{eq:scaling} gives the local
energy identity
\begin{equation}\label{eq:energy-scaling}
 \int_{B_1}(f_Ru_R+g_Rv_R)
 =R^{\alpha+\beta+2-n}\int_{B_R}(fu+gv).
\end{equation}
We shall obtain a universal upper bound for the left side, valid for
every scaled solution.  Positivity of the exponent in
\eqref{eq:energy-scaling} will then yield the final contradiction.
